\section{Conclusion}

$\eta_{\max}$ is a significant predictor at $K{=}3$ (pooled $\hat{\beta}_{\mathrm{std}}{=}0.243$) but becomes statistically non-significant at $K{\geq}5$ ($0.091$, $0.086$; CI crosses zero), while $K_{\mathrm{eff}}$'s pooled coefficient remains positive ($0.286$--$0.353$, CI excluding zero) across all panel sizes. Bayesian calibration nearly neutralizes composition effects ($R^2{\leq}0.17$). The diversity--capability confound ($r{=}0.70$--$0.81$) precludes causal attribution; all findings are associational.

For practitioners designing multi-judge panels: (1)~select judges by individual robustness ($\min \eta_{\max}$) rather than family diversity, as this yields the lowest ASR in all 12 conditions; (2)~use calibrated aggregation to reduce composition sensitivity, particularly at small panel sizes; (3)~monitor clean-condition $K_{\mathrm{eff}}$ as a deployable diversity metric ($\rho{\geq}0.91$ stability under 4/6 attacks). The $K{=}3$ to $K{=}5$ transition represents a practical design threshold: below it, a single weak judge dominates panel vulnerability; above it, panel composition effects matter more than individual weakness.

The positive diversity--vulnerability association instantiates the classical insight that diversity benefits require sufficiently capable ensemble members \citep{brown2005managing}, extended here to adversarial LLM evaluation. Separating diversity from capability (through training-based manipulation or capability-matched panel construction) remains the central open question. Extending to open-ended generation, increasing template diversity beyond 1--2 variants per attack, and testing under partial-knowledge adversaries would further test the generality of these associations.
